\documentclass[12pt]{article}
\usepackage[margin=1in]{geometry}
\usepackage{enumitem}
\usepackage{titlesec}
\usepackage{parskip}
\usepackage{graphicx}
\usepackage{xcolor}
\usepackage{hyperref}
\usepackage{fontenc}

\title{\textbf{Carey and Shugart 1995 RG Notes}\\
\large Carey, John M, and Matthew Soberg Shugart. ``Incentives to cultivate a personal vote: A rank ordering of electoral formulas'' [1995].\\
\textit{Electoral Studies} 14(4): 417-439.}
\author{}
\date{}

\begin{document}

\maketitle

Heavy connection to another paper on the list about legislators cultivating personal followings -- need to find which one

Core question: when does a candidate campaign primarily based on their own reputation versus their party's general reputation?

Four salient dimensions:

1) ``Ballot'' -- the extent to which leaders do (or do not) have control over the ballots/party lists, such that voters can (or cannot) disturb that list through their votes

2) ``Pooling'' -- ``Whether votes cast for one candidate of a given party also contribute to the number of seats won in the district by the party as a whole.'' That is, to what extent are votes for individuals counted as votes for the party whom they represent? High pooling might reflect SNTV and primary elections

3) ``Votes'' -- Do voters select a party, a candidate, or some series of candidates (i.e., both)?

See table 1 for the actual ranking -- closed list is most party-forward while SNTV and personal-list PR are most personal-forward

General takeaway seems to be that the more voters are able to express specific preferences through their ballot (rather than casting a general preference overall), the more valuable it becomes to campaign off of personal reputation / the more personal reputation matters than party reputation

An increase in district magnitude results in an increase in the value of personal reputation for every system OTHER THAN CLOSED LIST

See Crisp et al 2004 -- significant empirical testing of many of these core theories, and largely supports the tie between closed party list and non-individualistic behavior

\end{document}
